% GENERATED FILE. Edit canonical lessons, then run npm run generate:books.
% canonical-source-hash: fnv1a64:775e2dd9558328f4
% canonical-lessons: ZH-C527-xiadan, ZH-C527-baozhuang, ZH-C527-huiyuan, ZH-C527-youhui, ZH-C527-gezi

\chapter{Ordering Online}
\label{ch:zh-ordering-online}
\hlchaptermodality{527}

\begin{chapteropening}
\textbf{By the end of this chapter:} \emph{I can talk about placing an order, packaging, being a member, special deals and checked patterns.}

Complete the last lesson of chapter 527: I can talk about placing an order, packaging, being a member, special deals and checked patterns.
\end{chapteropening}

\section[xiàdān]{xiàdān — to place an order}

\label{lesson:ZH-C527-xiadan}



\begin{quote}
Before the new one: say the Mandarin for in large amounts, then the Mandarin for a small amount of.
\end{quote}

\subsection*{You'll want to know: xiàdān}
\emph{xiàdān} — \textquotedblleft{}to place an order\textquotedblright{}.

\begin{quote}
\emph{wǒ xiàdān yǐjīng sān tiān le dōngxi hái méi dào} — I placed the order three days ago, and it still hasn't arrived.
\end{quote}

Its characters come later; for now it is a word for the ear.

\begin{cousinweb}
\emph{xiàdān} is built from \zh{下} 'to put down' and \emph{dān} 'a slip of paper', both inherited from Middle Chinese, \emph{haeX} and \emph{tan} (Guangyun, in Baxter's transcription), which glosses the first 'to go down' and the second 'single'; the sense 'an order slip' is modern.
\end{cousinweb}

\begin{grammarlens}[title={What you've built}]
One more word for saying how it was.
\end{grammarlens}

\subsection*{Your turn}
\begin{itemize}
\raggedright
  \item \emph{Say it:} \emph{xiàdān}
  \item \emph{Say it:} \emph{xiàdān}, once more
  \item \emph{Say it:} \emph{wǒ xiàdān yǐjīng sān tiān le dōngxi hái méi dào} — I placed the order three days ago, and it still hasn't arrived.
  \item \emph{Recall:} say the Mandarin for in large amounts, then the Mandarin for a small amount of, then say \emph{xiàdān} again
\end{itemize}

\begin{tcolorbox}[breakable,colback=teal!4,colframe=teal!35!black,title={Before you move on}]
What does xiàdān mean? (\textquotedblleft{}To place an order\textquotedblright{}.) Say it once more.
\end{tcolorbox}
\section[bāozhuāng]{bāozhuāng — to wrap goods for sale}

\label{lesson:ZH-C527-baozhuang}



\begin{quote}
Before the new one: say the Mandarin for a small amount of, then the Mandarin for to place an order.
\end{quote}

\subsection*{You'll want to know: bāozhuāng}
\emph{bāozhuāng} — \textquotedblleft{}to wrap goods for sale\textquotedblright{}.

\begin{quote}
\emph{zhè jiā shāngdiàn de dōngxi bāozhuāng de zuì hǎokàn} — This shop packages its things the most beautifully.
\end{quote}

Its characters come later; for now it is a word for the ear.

\begin{cousinweb}
\emph{bāozhuāng} is built from \zh{包} 'to wrap' and \emph{zhuāng} 'to pack', both inherited from Middle Chinese, \emph{paew} and \emph{tsrjang} (Guangyun, in Baxter's transcription), which glosses the first 'to wrap up' and the second 'to bind up, to dress'.
\end{cousinweb}

\begin{grammarlens}[title={What you've built}]
One more word for saying how it was.
\end{grammarlens}

\subsection*{Your turn}
\begin{itemize}
\raggedright
  \item \emph{Say it:} \emph{bāozhuāng}
  \item \emph{Say it:} \emph{bāozhuāng}, once more
  \item \emph{Say it:} \emph{zhè jiā shāngdiàn de dōngxi bāozhuāng de zuì hǎokàn} — This shop packages its things the most beautifully.
  \item \emph{Recall:} say the Mandarin for a small amount of, then the Mandarin for to place an order, then say \emph{bāozhuāng} again
\end{itemize}

\begin{tcolorbox}[breakable,colback=teal!4,colframe=teal!35!black,title={Before you move on}]
What does bāozhuāng mean? (\textquotedblleft{}To wrap goods for sale\textquotedblright{}.) Say it once more.
\end{tcolorbox}
\section[huìyuán]{huìyuán — a member of a club or a shop's scheme}

\label{lesson:ZH-C527-huiyuan}



\begin{quote}
Before the new one: say the Mandarin for to place an order, then the Mandarin for to wrap goods for sale.
\end{quote}

\subsection*{You'll want to know: huìyuán}
\emph{huìyuán} — \textquotedblleft{}a member of a club or a shop's scheme\textquotedblright{}.

\begin{quote}
\emph{wǒ dāng le sān nián huìyuán} — I've been a member for three years.
\end{quote}

Its characters come later; for now it is a word for the ear.

\begin{cousinweb}
\emph{huìyuán} is built from \zh{会} 'to meet' and \emph{yuán}, both inherited from Middle Chinese, \emph{hwajH} and \emph{hjwen} (Guangyun, in Baxter's transcription), which glosses the first 'to come together' and the second as the number of things; the sense 'a member' is not given there.
\end{cousinweb}

\begin{grammarlens}[title={What you've built}]
One more word for saying how it was.
\end{grammarlens}

\subsection*{Your turn}
\begin{itemize}
\raggedright
  \item \emph{Say it:} \emph{huìyuán}
  \item \emph{Say it:} \emph{huìyuán}, once more
  \item \emph{Say it:} \emph{wǒ dāng le sān nián huìyuán} — I've been a member for three years.
  \item \emph{Recall:} say the Mandarin for to place an order, then the Mandarin for to wrap goods for sale, then say \emph{huìyuán} again
\end{itemize}

\begin{tcolorbox}[breakable,colback=teal!4,colframe=teal!35!black,title={Before you move on}]
What does huìyuán mean? (\textquotedblleft{}A member of a club or a shop's scheme\textquotedblright{}.) Say it once more.
\end{tcolorbox}
\section[yōuhuì]{yōuhuì — a favourable deal (for customers)}

\label{lesson:ZH-C527-youhui}



\begin{quote}
Before the new one: say the Mandarin for to wrap goods for sale, then the Mandarin for a member of a club.
\end{quote}

\subsection*{You'll want to know: yōuhuì}
\emph{yōuhuì} — \textquotedblleft{}a favourable deal (for customers)\textquotedblright{}.

\begin{quote}
\emph{huìyuán yǒu yōuhuì} — Members get a special deal.
\end{quote}

Its characters come later; for now it is a word for the ear.

\begin{cousinweb}
\emph{yōuhuì} is built from \emph{yōu} and \emph{huì}, both inherited from Middle Chinese, \emph{'juw} and \emph{hwejH} (Guangyun, in Baxter's transcription), which glosses the first 'plentiful' and the second 'kindness'.
\end{cousinweb}

\begin{grammarlens}[title={What you've built}]
One more word for saying how it was.
\end{grammarlens}

\subsection*{Your turn}
\begin{itemize}
\raggedright
  \item \emph{Say it:} \emph{yōuhuì}
  \item \emph{Say it:} \emph{yōuhuì}, once more
  \item \emph{Say it:} \emph{huìyuán yǒu yōuhuì} — Members get a special deal.
  \item \emph{Recall:} say the Mandarin for to wrap goods for sale, then the Mandarin for a member of a club, then say \emph{yōuhuì} again
\end{itemize}

\begin{tcolorbox}[breakable,colback=teal!4,colframe=teal!35!black,title={Before you move on}]
What does yōuhuì mean? (\textquotedblleft{}A favourable deal (for customers)\textquotedblright{}.) Say it once more.
\end{tcolorbox}
\section[\texorpdfstring{\zh{格子} (gézi)}{格子 (gézi)}]{\zh{格子} (gézi) — a checked pattern; squares}

\label{lesson:ZH-C527-gezi}



\begin{quote}
Before the new one: say the Mandarin for a member of a club, then the Mandarin for a favourable deal.
\end{quote}

\subsection*{You'll want to know: \zh{格子}}
\textbf{\zh{格子}} — \emph{gézi} — \textquotedblleft{}a checked pattern; squares\textquotedblright{}.

\begin{quote}
\emph{wǒ zuì xǐhuan zhè tiáo gézi qúnzi} — This checked skirt is my favourite.
\end{quote}

\begin{cousinweb}
\textbf{\zh{格子}} (\emph{gézi}) is built from \zh{格} 'a pattern' and \zh{子}, both inherited from Middle Chinese, \emph{kaek} and \emph{tsiX} (Guangyun, in Baxter's transcription); the second is said light, as at the end of many names of things.
\end{cousinweb}

\begin{grammarlens}[title={What you've built}]
One more word for saying how it was.
\end{grammarlens}

\subsection*{Your turn}
\begin{itemize}
\raggedright
  \item \emph{Say it:} \emph{gézi}
  \item \emph{Say it:} \emph{gézi}, once more
  \item \emph{Say it:} \emph{wǒ zuì xǐhuan zhè tiáo gézi qúnzi} — This checked skirt is my favourite.
  \item \emph{Recall:} say the Mandarin for a member of a club, then the Mandarin for a favourable deal, then say \emph{gézi} again
\end{itemize}

\begin{tcolorbox}[breakable,colback=teal!4,colframe=teal!35!black,title={Before you move on}]
What does \zh{格子} mean? (\textquotedblleft{}A checked pattern; squares\textquotedblright{}.) Say it once more.
\end{tcolorbox}
